\Th Пусть $F$ "--- конечное поле, $f \in F[x]$ "--- неприводимый многочлен. Тогда $f$ не имеет кратных корней над $\overline{F}$.

\Proof 
Пусть это не так, тогда $(f, f') \not\sim 1$. Тогда, в силу неприводимости многочлена $f$, $f' = 0$. Значит, $f$ имеет вид $f = a_mx^{mp} + \dotsc + a_1x^p + a_0$ для некоторых $a_0, \dotsc, a_m \in F$. Но тогда $f = (a_mx^{m} + \dotsc + a_1x + a_0)^p$, что противоречит его неприводимости.
\EndProof

\Th $\Aut\F_{p^n} \cong \Z_n$.

\Proof
Обозначим через $\phi \in \Aut\F_{p^n}$ \textit{автоморфизм Фробениуса}, имеющий вид $\phi(x) := x^p$ для всех $x \in \F_{p^n}$. Очевидно, $\phi$ сохраняет операции, инъективность же выполнена потому, что $\ke\phi = \{0\}$, и ее достаточно для биективности, поскольку $\F_{p^n}$ конечно. Заметим также, что $\F_{p^n}^\phi = \{a \in \F_{p^n} : a^p = a\} = \Z_p$.
	
	С одной стороны, $\phi^n = \id$. С другой стороны, $\forall k \in \{1, 
	\dotsc, n - 1\}:$ равенство $\phi^k(x) = x$ выполнено только для таких элементов поля $\F_{p^n}$, которые являются корнями многочлена $x^{p^k} - x \in \Z_p[x]$.  Но $|\Aut\F_{p^n}| = |\Aut_{\Z_p}\F_{p^n}| = [\F_{p^n} : \Z_p] = n$ в силу нормальности расширения $\F_{p^n} \supset \Z_p$, поэтому $\Aut\F_{p^n} = \gl\phi\gr = \{\id, \phi, \dotsc, \phi^{n-1}\} \cong \Z_n$.

\EndProof 